\pdfoutput=1
\documentclass[11pt]{article}
\usepackage[margin=1in]{geometry}
\usepackage{amsmath,amssymb,amsthm}
\usepackage{enumitem}
\usepackage{booktabs}
\usepackage{hyperref}
\hypersetup{hidelinks,
  pdftitle={The cost of each distinction in the local computations of IUT IV},
  pdfauthor={Toshihisa Urahigashi}}

\theoremstyle{plain}
\newtheorem{theorem}{Theorem}
\newtheorem{proposition}[theorem]{Proposition}
\newtheorem{corollary}[theorem]{Corollary}
\theoremstyle{definition}
\newtheorem{remark}{Remark}

\newcommand{\Fmod}{F_{\mathrm{mod}}}
\newcommand{\ord}{\operatorname{ord}}
\newcommand{\Norm}{\operatorname{N}}
\newcommand{\mulog}{\mu^{\log}}

\title{The cost of each distinction in the local computations\\
of inter-universal Teichm\"uller theory IV}
\author{Toshihisa Urahigashi\\
\small\texttt{123026.urahigashi.toshihisa@gmail.com}}
\date{Version 1.0.0 --- 7 September 2026\\[3pt]
\small ORCID \href{https://orcid.org/0009-0004-0460-6242}{\texttt{0009-0004-0460-6242}}\\[3pt]
\small DOI \href{https://doi.org/10.5281/zenodo.22646137}{10.5281/zenodo.22646137}}

\begin{document}
\maketitle

\begin{abstract}
The local estimates of \cite[\S1]{IUT4} are assembled from several steps that are kept
apart in the source: a $p$-primary torsion term, a normalization that divides by the
local degree, a tensor order sitting inside its own normalization, a floor function, a
ramification support, and a weighting of places.  Each separation is drawn explicitly
there.  This report computes, in actual local fields and for named curves, \emph{what
each of them is worth}: how much the estimate moves if that one distinction is elided
and nothing else is changed.  The answers are exact rationals or exact multiples of
$\log p$ --- $\log3/2$, $\log5/4$, $\log5/4$ again for a different reason, $156/157$,
a factor of $5.49\times10^{4}$, a weight of $2355/2356$ against $1/2$, and $\log661$ at
a single prime.  \path{scripts/costs.py} computes twelve tables --- standard library
only, each asserting its own conclusion, exit zero --- of which the ones displayed here
are those the argument uses.

\textup{[S]} marks a formula or hypothesis printed in the source and \textup{[P]} an
exact computation from it.  \emph{Nothing here is a correction to \cite{IUT4}.} Every
distinction measured below is one the source itself makes --- with the single
exception of \S\ref{s:wt}, noted there; the numbers say what an
implementation loses by collapsing it, which is a question only an explicit datum can
answer.  \emph{No claim is made about \cite[Cor.~3.12]{IUT3}, about the objections of
\cite{SS}, or about $abc$.}
\end{abstract}

\section{What this is}\label{s:what}

\textup{[S]} \cite[Thm.~1.10]{IUT4} fixes a collection of initial $\Theta$-data as in
\cite[Def.~3.1]{IUT1}, but \emph{it is not stated for an arbitrary one}: it supposes
in addition that one is in the situation of \cite[Cor.~3.12]{IUT3}, and that $E_F$ has
good reduction at every $v\in\mathbb V(F)^{\mathrm{good}}\cap\mathbb V(F)^{\mathrm{non}}$
not dividing $2\ell$.  \textup{[P]} Its proof passes through a chain of local statements --- \cite[Prop.~1.1]{IUT4}
through \cite[Prop.~1.7]{IUT4} --- each of which separates quantities that a reader
might be tempted to identify.  This report evaluates them on named local fields and
named curves.  \emph{Several of them are computable without a full initial
$\Theta$-datum} --- and \emph{some of the curves used below are not known to carry
one}, the $N=472$ curve of \cite[\S13]{Second} being a curve with fourteen
\emph{candidate} levels.  What the verified data of \cite{First,Second} add is
\S\ref{s:bdry} and \S\ref{s:recovered}, where the parameters of an actual datum enter.

\textup{[P]} The pattern is the same in every section.  We take the source's own
statement, instantiate it in a local field or at a curve that is named, compute the
quantity exactly, then compute it again with one distinction removed, and print the
difference.  The difference is the cost of that distinction.  It is a number about an
implementation, not about the theory: \emph{a bound that keeps all the distinctions is
correct whether or not the cost of any one of them is large.}

\begin{remark}[why this is worth writing down]\label{r:why}
\textup{[P]} Two of the costs below were found because an earlier draft of
\cite{Second} got them wrong --- an incorrect ramification index, and a formula applied
where its hypothesis failed.  In both cases what caught the error was the source's own
statement being sharp enough to check against a specific curve.  A table of costs is
the form in which that sharpness becomes reusable.
\end{remark}

\section{Notation}\label{s:not}

\textup{[S]} For a finite extension $k/\mathbb{Q}_p$ with ring of integers
$\mathcal{O}$, ramification index $e$ and residue degree $f$, we write $\mulog$ for the
log-volume normalized by $[k:\mathbb{Q}_p]$, as in \cite[Prop.~1.4(i)]{IUT4} --- the
notation of \cite[Prop.~1.1, Prop.~1.2]{IUT4} supplies $d_i$ and $a_i$, which enter
later --- so that
$\mulog(\mathcal{O})=0$ and $\mulog$ of a sublattice of index $p^{m}$ is
$-\,m\log p/[k:\mathbb{Q}_p]$.  All values below are given in units of $\log p$ unless
stated otherwise.

\section{The $p$-primary torsion term}\label{s:tors}

\textup{[S]} \cite[Prop.~1.4(ii)]{IUT4} gives, for a finite extension of
$\mathbb{Q}_p$ with ramification index $e$, residue degree $f$, and $p$-primary torsion
of order $p^{m}$ in the units,
\[
 \mulog\bigl(\log\mathcal{O}^{\times}\bigr)=-\Bigl(\frac1e+\frac{m}{ef}\Bigr)\log p .
\]
\textup{[P]} Take $k=\mathbb{Q}_p(\zeta_p)$, so $e=p-1$, $f=1$ and $m=1$.  Dropping the
torsion term deletes $m/(ef)$:

\begin{center}
\begin{tabular}{rrrrrrr}
\toprule
$p$ & $e$ & $f$ & $m$ & $\mulog/\log p$ & without $m$ & cost\\
\midrule
3 & 2 & 1 & 1 & $-1$ & $-1/2$ & $1/2$\\
5 & 4 & 1 & 1 & $-1/2$ & $-1/4$ & $1/4$\\
7 & 6 & 1 & 1 & $-1/3$ & $-1/6$ & $1/6$\\
\bottomrule
\end{tabular}
\end{center}

\textup{[P]} So the cost is $\tfrac12\log3$, $\tfrac14\log5$, $\tfrac16\log7$: it is
\emph{not} uniform in $p$, and it is largest where $p$ is smallest.

\begin{remark}[correction]\label{r:tors}
\textup{[P]} An earlier version derived these from the ceiling in $a_i$ of
\cite[Prop.~1.2]{IUT4}.  The values were right and the derivation was not: $a_i$ is a
coefficient of the hull bound and $m_i$ is the torsion variable, and they are different
objects.  The normalization itself is \cite[Prop.~1.4(i)]{IUT4}.
\end{remark}

\section{The normalization by the local degree}\label{s:norm}

\textup{[S]} \cite[Prop.~1.4(i),(ii)]{IUT4} states its bounds for the normalized volume.
\textup{[P]} In the unramified extensions of $\mathbb{Q}_5$ of degrees $1,2,3$ the
lattice index moves and the normalized volume does not:

\begin{center}
\begin{tabular}{rrrr}
\toprule
$f$ & index & raw $\log$ index$/\log5$ & $\mulog/\log5$\\
\midrule
1 & 5 & 1 & $-1$\\
2 & 25 & 2 & $-1$\\
3 & 125 & 3 & $-1$\\
\bottomrule
\end{tabular}
\end{center}

\textup{[P]} Reading the raw index as a log-volume therefore counts the residue degree
a second time, and the error grows linearly in $f$.  This is the cheapest distinction to
state and the easiest to lose, because both quantities are ``the index''.

\section{The tensor order inside its normalization}\label{s:tens}

\textup{[S]} \cite[Prop.~1.1]{IUT4} and \cite[Prop.~1.4(iii)]{IUT4} distinguish an order
from the normalization in which it sits.  \textup{[P]} Let $k=\mathbb{Q}_5(\sqrt5)$,
$\pi=\sqrt5$, and $R=\mathcal{O}\otimes_{\mathbb{Z}_5}\mathcal{O}$ inside
$\mathcal{O}\times\mathcal{O}$.  Then
\[
  R=\{(u,v)\in\mathcal{O}^{2}: u\equiv v \bmod \pi\},\qquad
  [\mathcal{O}^{2}:R]=|\mathcal{O}/\pi|=5,\qquad
  \dim_{\mathbb{Q}_5}(k\times k)=4,
\]
in which $\mathcal{O}^{2}$ is a $\mathbb{Z}_5$-lattice of rank $4$ and $k\times k$ is
the $\mathbb{Q}_5$-vector space it spans,
so $\mulog(R)=-\tfrac14\log5$.  \textup{[P]} Treating $R$ as its own normalization sets
this to $0$ and loses exactly $\log5/4$.  \emph{Numerically this is the same as the cost of
the torsion term at $p=5$ in \S\ref{s:tors}, namely $1/(p-1)=1/4$, and it has a
different origin;} the coincidence is a reason to keep the two apart rather than to
merge them.  \emph{It is not the same as anything in \S\ref{s:norm}}: there the $f=4$
discrepancy between the raw and the normalized volume would be $3\log5$.

\section{The floor}\label{s:floor}

\textup{[S]} \cite[Prop.~1.4(iii)]{IUT4} passes from the container to the smooth bound
by $-\lfloor A\rfloor\le-A+1$, so the slack is $1+\lfloor A\rfloor-A$.  \textup{[P]} In
the moderate tame model here $d_I+a_I=|I|$ is an integer, so the slack is
$1+\lfloor\lambda\rfloor-\lambda$:

\begin{center}
\begin{tabular}{rrr}
\toprule
$\lambda$ & $1+\lfloor\lambda\rfloor-\lambda$ & $\lceil\lambda\rceil-\lambda$ (wrong)\\
\midrule
$3$ & $\mathbf 1$ & $0$\\
$472/157$ & $156/157$ & $156/157$\\
$627/157$ & $1/157$ & $1/157$\\
$4$ & $\mathbf 1$ & $0$\\
\bottomrule
\end{tabular}
\end{center}

\begin{remark}[correction]\label{r:floor}
\textup{[P]} An earlier version used $\lceil\lambda\rceil-\lambda$.  The two agree off
the integers and differ \emph{on} them: at $\lambda=3$ and $\lambda=4$ the slack is $1$,
not $0$.
\end{remark}

\textup{[P]} At $\lambda=472/157$ --- the value carried by the curve of
\cite[\S13]{Second}, which has \emph{fourteen candidate levels} --- the slack is
$156/157$, while the stage-one entry margin of that same curve is $1/157$.  \emph{Their ratio is $156$, and they are not
comparable quantities at all:} one is a defect between a container and a smooth bound, the other a margin
between the two sides of a comparison.  Placing them side by side without saying which
is which is easy to lose, and that is why the two are printed together here.

\section{The ramification support of the input}\label{s:supp}

\textup{[S]} \cite[Thm.~1.10]{IUT4}, Step (iii), and \cite[Rem.~1.10.5(iii)]{IUT4} work
with the ramification support of the input.  \emph{Step (viii) then bounds that support
from above by a sum over all primes below a cutoff}, so the source draws both, and what
is measured here is the size of the gap between them at one curve.  An earlier version
of this paragraph said the source used the support \emph{rather than} all primes; that
misdescribed Step (viii) (external review, 2026-09-07).  \textup{[P]} Take the curve
\texttt{5190c3}, with
\[
 \Delta = 2^{6}\cdot3^{44}\cdot5\cdot173 ,
\]
and the torsion field $K_0=\mathbb{Q}(i,E[210])$.  Its ramification support is
$\{2,3,5,7,173\}$ --- five primes.  \emph{That this list is exactly the support is an
ordinary number-theoretic statement about the curve and the torsion field, not a
consequence of the factorization}.  It is proved as follows.  Away from those five
primes $E$ has good reduction and $210$ is invertible, so by the criterion of
N\'eron--Ogg--Shafarevich the field $K_0$ is unramified there.  Conversely $2$ ramifies
in $\mathbb{Q}(i)$; the Weil pairing puts $\mathbb{Q}(\mu_{210})$ inside $K_0$, and
$3,5,7$ ramify in it; and at $173$ one has $v(\Delta)=1$ with $c_4$ a unit, so the
reduction is multiplicative and the Tate parameter has valuation $1$ after at most an
unramified quadratic extension --- rationalizing the $3$-torsion then requires an
element of valuation $1/3$, so $173$ ramifies as well.  Both inclusions follow.
\textup{[P]} What the script checks is the arithmetic this rests on: what the script checks is that it is
consistent with $\Delta$ and with the ramification of $\mathbb{Q}(i,\mu_{210})$.  The number of \emph{all} primes below the cutoff
$3\,870\,720$ is $274\,666$, computed by sieve, with no appeal to any effective prime
number theorem.  The ratio is
\[
  \frac{274666}{5}=54933.2 .
\]

\textup{[P]} So replacing the input support by the all-prime proxy that Step (viii)
uses multiplies the count of primes by $5.49\times10^{4}$ at this curve.  \emph{It is the only cost
tabulated here that is a factor rather than an additive term}; the comparison is of
counts of primes, not of every quantity in the report.

\section{The weighting of places}\label{s:wt}

\textup{[S]} \cite[Rem.~1.7.1]{IUT4} weights the places of the moduli field by their
local degrees over $[\Fmod:\mathbb{Q}]$.  \textup{[P]} \emph{Weighting them by their
ramification indices instead is not an alternative the source anywhere draws}: unlike
the other sections here, this one measures an elision that only an implementation could
make.  At a split prime with $e_{\mathrm b}=2355$ and $e_{\mathrm g}=1$:

\begin{center}
\begin{tabular}{lr}
\toprule
weight from $e$ alone, bad side & $2355/2356$\\
correct weight, bad side & $1/2$\\
ratio & $2355/1178$\\
\midrule
two-factor tuples $\mathrm{gg},\mathrm{gb},\mathrm{bg},\mathrm{bb}$ & $1/4$ each\\
\bottomrule
\end{tabular}
\end{center}

\textup{[P]} Weighting by $e$ sends the bad side to $2355/2356$, that is, to almost all
of the mass, where the correct weight gives it half.  The auxiliary degrees cancel; the
place count does not.

\section{A mean that agrees while the components do not}\label{s:mean}

\textup{[S]} \cite[Prop.~1.7]{IUT4} sums over ordered tuples, and
\cite[Rem.~1.10.7(ii)]{IUT4} keeps the mixed components.  \textup{[P]} In
$(\mathbb{Q}_p\times\mathbb{Q}_p)\otimes(\mathbb{Q}_p\times\mathbb{Q}_p)$ there are four
components.  \emph{Assigning them the illustrative values $0,8,8,16$} --- these are
chosen to display the phenomenon, not read off a datum --- the mean is $8$, and the two
diagonal components alone, with $0,16$, have mean $8$ as well.

\textup{[P]} \emph{The means agree.} What differs is the probability carried by a mixed
component: $1/2$ against $0$.  So a check on the mean cannot detect the collapse of the
mixed components, and any verification of an implementation that compares only averages
will pass a program that has lost them.

\section{Two boundaries in the explicit families}\label{s:bdry}

\textup{[P]} $e_{\mathrm b}\le p-2$ is \emph{not} a hypothesis of
\cite[Prop.~1.4(iii)]{IUT4}: there $I^{*}$ is a chosen subset and the inequality is
asked only outside it.  It is the condition under which the place may be left out of
$I^{*}$, that is, the condition for the error-free form of the bound.  \textup{[P]}
For $\pi=41+26i$, $p=2357$, $\ell=157$ and $a=\pi^{N}$ one has
$e_{\mathrm b}=15\ell/\gcd(15\ell,2N)$, so a single step in $N$ moves it by a factor of
$471$:

\begin{center}
\begin{tabular}{rrrr}
\toprule
$N$ & $\gcd(2355,2N)$ & $e_{\mathrm b}$ & $e_{\mathrm b}\le2355$\\
\midrule
471 & 471 & 5 & yes\textsuperscript{$\dagger$}\\
\textbf{472} & \textbf{1} & \textbf{2355} & yes\\
474 & 3 & 785 & yes\\
475 & 5 & 471 & yes\\
\bottomrule
\end{tabular}
\end{center}

\textup{[P]} \textsuperscript{$\dagger$}\emph{The row $N=471$ is a local arithmetic
candidate only}: there $\ell=157$ divides $2N=942$, so \cite[Def.~3.1(c)]{IUT1} is not
satisfied and it is not a datum to compare with.  The row is printed because the gcd
computation is the point.  The curve of \cite[\S13]{Second} --- the one with
\emph{fourteen candidate levels}, not fourteen verified initial $\Theta$-data --- has
$N=472$, where the gcd is $1$ and $e_{\mathrm b}$ is maximal.

\textup{[P]} The second boundary concerns gcds of ideals against gcds of norms.  For
$\pi=16+3\sqrt5$ and $N=110$ the prime $661$ splits, and $\sqrt5$ has the two roots
$115$ and $546$ modulo $661$.  \emph{Those residues are not square roots of $5$ to
higher precision} --- $115^{2}-5=20\cdot661$ is not $0$ modulo $661^{2}$ --- so they are
Hensel-lifted before use; deciding that a valuation is exactly $1$ needs a genuine
local square root, not a root modulo $661$.  Reducing $a=(16+3\sqrt5)^{110}$ and
$t=1-a+a^{2}$ at each lifted root:

\begin{center}
\begin{tabular}{rrrr}
\toprule
$\sqrt5\bmod661$ & $a\bmod661$ & $v(1-a)$ & $v(t)$\\
\midrule
115 & 1 & 1 & 0\\
546 & 365 & 0 & 1\\
\bottomrule
\end{tabular}
\end{center}

\textup{[P]} \emph{The $661$ of the numerator and that of the denominator sit at
different primes.}  The total norms nonetheless carry $661$ twice, so cancelling them
removes it once too often.  Writing $D=\Norm(a^{2}b^{2})$ and $U=\Norm(256t^{3})$ for
$b=1-a$, and taking the Legendre invariant $j=256t^{3}/(a^{2}b^{2})$, the two finite
heights normalized by the degree are
\[
  H_{\mathrm{ideal}}=\tfrac12\log\frac{D}{G_{\mathrm{ideal}}},\qquad
  H_{\mathrm{norms}}=\tfrac12\log\frac{D}{\gcd(U,D)} ,
\]
where $G_{\mathrm{ideal}}=\Norm\bigl((256t^{3},a^{2}b^{2})\bigr)=2^{8}$.
\emph{The identity $t+ab=1$ fixes the support of the common factor, not its exponent},
so the $2$-adic valuation has to be computed.  Since $t$ and $ab$ share no prime ideal,
the only common prime is the $2$ coming from $256$; and $2$ is inert in
$\mathbb{Q}(\sqrt5)$, so $v_{(2)}(x)=\tfrac12 v_2(\Norm x)$.  Here $a\equiv5$ and
$b\equiv4\pmod 8$, whence $v_{(2)}(a)=v_{(2)}(t)=0$ and $v_{(2)}(b)=2$.  Therefore
$v_{(2)}(256t^{3})=8$ and $v_{(2)}(a^{2}b^{2})=4$, so the ideal gcd is $(2)^{4}$ and its
norm is $2^{8}$.  The norm gcd carries $661^{2}$ in
addition.  Both are computed from the curve in \path{scripts/costs.py}, giving
\[
  H_{\mathrm{ideal}}=1174.636200642495,\qquad
  H_{\mathrm{norms}}=1168.142446802643,\qquad
  H_{\mathrm{ideal}}-H_{\mathrm{norms}}=\log661
\]
to twelve places.  \emph{The reason is in the table, not only in the decimals}: a
check carried out on the norms alone cannot see which prime a valuation
belongs to.

\section{A finite statistic for the tame local coefficient}\label{s:statistic}

\textup{[P]} \emph{This section is about the coefficient of $\log p$ in the simplified
bound, not about an actual possible-image log-volume.}  Fix a stage cutoff $m\ge2$ and
write
\[
 T_m=\coprod_{j=1}^{m}\{\mathrm g,\mathrm b\}^{j+1},\qquad
 c_{\mathrm g}=1-e_{\mathrm g}^{-1},\quad c_{\mathrm b}=1-e_{\mathrm b}^{-1},\quad
 s=\ord_p(\underline{\underline q})>0 .
\]
\textup{[P]} In the tame specialization with every $e_i\le p-2$ --- so that
$I^{*}=\varnothing$ is available, which is a \emph{specialization} of
\cite[Prop.~1.4(iii)]{IUT4} and not a hypothesis satisfied by arbitrary tuples --- the
coefficient collapses to two counts and one indicator:
\[
 \boxed{\;B(t)=1+c_{\mathrm g}n_{\mathrm g}(t)+c_{\mathrm b}n_{\mathrm b}(t)
        -s\,j(t)^2\mathbf 1_{\{t_j=\mathrm b\}}\;}
\]

\begin{proposition}\label{p:statistic}
\textup{[P]} Put $\pi(t)=(n_{\mathrm b}(t),\bot)$ when $t_j=\mathrm g$ and
$\pi(t)=(n_{\mathrm b}(t),j)$ when $t_j=\mathrm b$.  \emph{On $T_m$} with $m\ge2$, $B$
factors through $\pi$ if and only if $e_{\mathrm g}=1$.  If moreover $e_{\mathrm b}>1$
and $s>2c_{\mathrm b}$, the induced map is injective.
\end{proposition}

\begin{proof}
Sufficiency is substitution.  For necessity, $(\mathrm g,\mathrm g)$ and
$(\mathrm g,\mathrm g,\mathrm g)$ have the same $\pi$ and their $B$ differ by
$c_{\mathrm g}$; this uses $m\ge2$.  For injectivity write $c=c_{\mathrm b}>0$: the
good-ending values are $\ge1$ while the largest bad-ending value at stage $1$ is
$1+2c-s<1$, and the maximum at stage $j+1$ minus the minimum at stage $j$ is
$c(j+1)-s(2j+1)<0$, so the stages occupy disjoint ranges.
\end{proof}

\textup{[P]} \emph{The domain is part of the statement.}  At a \emph{fixed} stage
$n_{\mathrm g}=j+1-n_{\mathrm b}$, so $B$ factors through $\pi$ for every
$e_{\mathrm g}$ and the factorization says nothing; the equivalence lives on $T_m$,
where the two counts are independent.  The strict inequality is also needed: at
$s=2c_{\mathrm b}$ the classes $(0,\bot)$ and $(2,1)$ take the same value.

\begin{remark}[the exact criterion, and why $s>2c_{\mathrm b}$ suffices]
\label{r:exact-criterion}
\textup{[P]} The condition $s>2c_{\mathrm b}$ is \emph{sufficient but not necessary}.
Put $r=s/c_{\mathrm b}$ and
\[
 E_m=\Bigl\{\tfrac{d}{j^{2}}:1\le j\le m,\ 1\le d\le j+1\Bigr\}\cup
      \Bigl\{\tfrac{d}{k^{2}-j^{2}}:1\le j<k\le m,\ 1\le d\le k\Bigr\}.
\]
Then, for $e_{\mathrm g}=1$ and $c_{\mathrm b}>0$, \emph{the induced map is injective
if and only if $r\notin E_m$}: a good-ending class meets a bad-ending one at stage $j$
exactly when $r=(n-h)/j^{2}$ with $1\le n-h\le j+1$, and two bad-ending stages $j<k$
meet exactly when $r=(n_k-n_j)/(k^{2}-j^{2})$ with $1\le n_k-n_j\le k$; within one
group $c_{\mathrm b}>0$ separates.  \emph{This criterion is due to an external
reviewer of the present note.}

\textup{[P]} It also explains the crude condition.  The first family is maximised at
$j=1$, $d=2$, giving $2$; the second is bounded by $k/(k^{2}-(k-1)^{2})=k/(2k-1)<1$.
Hence $\sup E_m=2$ \emph{for every} $m$, so $r>2$ --- that is, $s>2c_{\mathrm b}$ ---
already misses $E_m$ at every cutoff.  For the profile of \cite{Second},
$r=\tfrac{29}{7}\big/\tfrac{104}{105}=\tfrac{435}{104}>2$.
\end{remark}

\textup{[P]} For the profile of \cite{Second} --- $e_{\mathrm g}=1$,
$e_{\mathrm b}=105$, $s=29/7$, so $c_{\mathrm b}=104/105$ and $2c_{\mathrm b}<s$ --- the
hypotheses hold, and $T_3$ has $28$ tuples with $13$ distinct values.  \emph{So the
good places are invisible to this coefficient}: it sees how many bad places there are,
and, when the last place is bad, at which stage it is applied.

\begin{remark}\label{r:statistic-scope}
\textup{[P]} This is a minimal statistic \emph{for the single function $B$}: its fibres
agree with those of $\pi$ under the stated hypotheses.  It is not a claim that the
statistic supports any later operation in the source's proof, and it is not a statement
about actual log-volumes.
\end{remark}

\section{Two items recovered}\label{s:recovered}

\textup{[P]} Two items were listed as not established in an earlier version of this
report.  Both are now settled, and in each case the author's own computation had been
correct while the quantity being compared was not the one reported.

\paragraph{The two-adic boundary.}
\textup{[P]} For $\pi=10+3i$ the finite height contribution at $2$ is not $v_2$ itself
but the remainder left by a cap at $8$.  \emph{The cap is ours, and it comes from this
family, not from a general statement}: \cite[Cor.~2.2(i)]{IUT4} defines
$\log(q^{\forall})$ and $\log(q^{\neq2})$ and prints no such cutoff.  At the prime
$\mathfrak{p}=(1+i)$ of $\mathbb{Q}(i)$ one has $v_{\mathfrak p}(2)=2$ and
$v_{\mathfrak p}(a)=v_{\mathfrak p}(1-a+a^2)=0$, so from $j=256t^3/(a^2b^2)$,
\[
 v_{\mathfrak p}(j)=16-2v_{\mathfrak p}(1-a),\qquad
 v_{\mathfrak p}(1-a)=v_2\bigl(\Norm(1-a)\bigr),
\]
and the normalized local finite height $\tfrac12\max\{-v_{\mathfrak p}(j),0\}\log2$ is
\[
 \max\bigl\{v_2(\Norm(1-a))-8,\ 0\bigr\}\log2 .
\]

\begin{center}
\begin{tabular}{rrrl}
\toprule
$N$ & $v_2(\Norm(1-a))$ & $\max\{v_2-8,0\}$ & contribution\\
\midrule
199 & 1 & 0 & $0$\\
200 & 9 & 1 & $\log2$\\
201 & 1 & 0 & $0$\\
400 & 11 & 3 & $3\log2$\\
\bottomrule
\end{tabular}
\end{center}

\textup{[P]} \emph{The values $1,9,1,11$ computed earlier were right}; they are $v_2$,
and the reported quantity is the excess over the cap.  One step in $N$ moves the
contribution between $0$ and $\log2$.

\paragraph{Two permitted upper bounds at one prime.}
\textup{[P]} At $\ell=157$, $p=2357$ one has $e_{\mathrm b}=2355$ and
$e_{\mathrm g}=1$, both $\le p-2=2355$, so $I^{*}=\varnothing$ is available and the
extra error term of \cite[Prop.~1.4(iii)]{IUT4} is \emph{exactly $0$}.  The uniform
coarse term, averaged over $j$, is
\[
 (\ell+5)\log(e^{*}_{\mathrm{mod}}\ell)=162\log(1105920\cdot157)
 =3073.534296832696\ldots
\]

\textup{[P]} \emph{This compares two permitted upper-bound expressions at one fixed
prime}, averaged over $j$.  It is not the actual indeterminacy, and it is not a gap in
\cite[Cor.~3.12]{IUT3}; no such claim is made here.

\section{What is not established here}\label{s:notdone}

\textup{[P]} \emph{Nothing remains on this list.}  The two items an earlier version
placed here --- the two-adic boundary and the comparison of a precise against a coarser
local error term --- are established in \S\ref{s:recovered}.  They are recorded here
because the reason for the earlier omission is worth keeping --- \emph{both are
resolved in \S\ref{s:recovered} of the present note and neither remains open} --- in
both cases the author's arithmetic was correct and the \emph{quantity being compared}
was not the one reported, so the check failed for a reason that no amount of
recomputation would have
found.

\section{Scope}\label{s:scope}

\noindent\textbf{Claimed.}  The values printed above, as exact rationals or exact
multiples of $\log p$, for the local fields and curves named with them.

\smallskip\noindent\textbf{Not claimed.}  That any statement of \cite{IUT1,IUT3,IUT4} is
in error; that any estimate in those papers is affected; that the $abc$ conjecture, the
objections of \cite{SS}, or \cite[Cor.~3.12]{IUT3} are touched by anything here.  No
log-volume of an actual $\Theta$-pilot object is computed, no hull is compared, and no
membership in any exceptional set is addressed.  \emph{The costs above are costs of
eliding a distinction in an implementation; they are not defects of the source, which
draws every one of those distinctions explicitly and is cited for each.}

\section*{Reproduction}

\path{scripts/costs.py} prints all twelve tables using only the standard library,
asserts each conclusion --- including the sieve count $\pi(3\,870\,720)=274666$, the
factorization of $\Delta$, the four values of $e_{\mathrm b}$, and the agreement of the
two means --- and exits zero.  \emph{Table 12 checks every clause of
Proposition~\ref{p:statistic}}: the equivalence with $e_{\mathrm g}=1$ on $T_3$, that
at a \emph{fixed} stage the factorization holds for every $e_{\mathrm g}$ and so would
be vacuous there, the injectivity under $e_{\mathrm b}>1$ and $s>2c_{\mathrm b}$, the
collision of $(0,\bot)$ with $(2,1)$ at $s=2c_{\mathrm b}$ exactly, and the counts
$|T_3|=28$ with $13$ distinct values.  \emph{That proposition is the one statement
here that was wrong twice before it was right} --- first in its domain, then in its
injectivity hypothesis --- which is why each clause is asserted rather than only
stated.

\textup{[P]} \emph{Four different kinds of warrant appear in this note and we do not
run them together.}  \emph{(1)} Exact rational or integer computation: the tables, and
the two heights of \S\ref{s:bdry}, computed from the curve.  \emph{(2)} Finite
exhaustive check over a stated range: Table 12's clauses, verified on $T_3$ and on the
rational $r$ enumerated there --- a check on a finite set, not a proof for all $m$.
\emph{(3)} Ordinary proof: Proposition~\ref{p:statistic} and
Remark~\ref{r:exact-criterion}, which are proved for all $m$ in the text and are
\emph{not} formalised anywhere.  \emph{(4)} Kernel verification, \emph{narrowly}: in
\cite{Second} the files \path{ImaginaryQuadraticDatum.lean} and
\path{ModSevenImage.lean} re-prove the integer identities and finite predicates of the
datum $a=(10+3i)^{29}$ --- norm factorisations, valuations, point counts and traces at
the witness places, and the finite group computations --- and
\path{FamilyArithmetic.lean} the finite algebra behind that report's existence family;
in \cite{First} the corresponding statements for $a=(16+3\sqrt5)^{29}$.  \emph{No
statement of the present note is among them}, and nothing here is discharged by them.
In particular the curve of \cite[\S13]{Second} with $N=472$ carries \emph{fourteen
candidate levels}, not fourteen verified initial $\Theta$-data.

\textup{[P]} Within (1), one more line is worth drawing: $D$, $U$, $G_{\mathrm{ideal}}$
and $\gcd(U,D)$ are exact integers and the displayed formulas for
$H_{\mathrm{ideal}}$, $H_{\mathrm{norms}}$ are exact; the twelve-place decimals printed
for them are floating-point evaluations of those exact expressions.

\section*{Disclosure}

Large language models were used substantially in locating passages, carrying out
computations, drafting and adversarial review; the author is responsible for the claims.
Every number above was recomputed by the author independently of the process that first
produced it, and the two items where an earlier recomputation had failed were
recorded rather than dropped silently; both are settled in \S\ref{s:recovered}.

\begin{thebibliography}{Second}
\bibitem[First]{First} T.~Urahigashi, \emph{An explicit initial $\Theta$-datum with
 mixed reduction at a split prime}, Zenodo, DOI 10.5281/zenodo.22537342, 2026.
\bibitem[Second]{Second} T.~Urahigashi, \emph{An initial $\Theta$-datum over an
 imaginary quadratic field, and numerical values of the local quantities of
 inter-universal Teichm\"uller theory IV}, 2026.
\bibitem[DH25]{DH25} T.~Dupuy, A.~Hilado, \emph{The statement of Mochizuki's
 Corollary 3.12: initial theta data}, arXiv:2004.13228v2 [math.NT], 19 June 2025.
\bibitem[IUT1]{IUT1} S.~Mochizuki, \emph{Inter-universal Teichm\"uller theory I},
 Publ. Res. Inst. Math. Sci. \textbf{57} (2021), 3--207.
\bibitem[IUT3]{IUT3} S.~Mochizuki, \emph{Inter-universal Teichm\"uller theory III},
 Publ. Res. Inst. Math. Sci. \textbf{57} (2021), 403--626.
\bibitem[IUT4]{IUT4} S.~Mochizuki, \emph{Inter-universal Teichm\"uller theory IV},
 Publ. Res. Inst. Math. Sci. \textbf{57} (2021), 627--723.
\bibitem[SS]{SS} P.~Scholze, J.~Stix, \emph{Why $abc$ is still a conjecture},
 manuscript, July 16, 2018.
\end{thebibliography}

\end{document}
